% 第一篇后半：PDF 物理页 30--54，原书第 25--49 页。
\origpage{30}{25}
\noindent b) 映射 $\CH^2(X)\to\CH^2(\overline X)$ 显然将前一群映入不变子群 $\CH^2(\overline X)^G$，而其核由转移论证可知是挠群。因此，要证明定理，只需证明
\[
 \Ker[\CH^2(X)\longrightarrow\CH^2(\overline X)]
\]
\term{指数有界}{finite exponent}{exposant fini}。事实上，对特征为零的域 $k$ 上任意光滑射影簇，只要 $H^1(X,\mathcal O_X)=0$ 且 $H^2(X,\mathcal O_X)=0$，这一结论就成立。[CR1] 已证明了这一点，下面简述其中的主要论证。\translatornoteref{2}

这里采用的是 S.~Bloch 在 [B4] 中首次使用、并由 [CR1] 系统化的\term{伽罗瓦方法}{Galois method}{méthode galoisienne}。其做法是：对 $\overline X$ 上层 $\mathcal K_2$ 的 Quillen 分解的整体截面复形取伽罗瓦上同调，由此得到正合列
\[
\begin{split}
\cdots\longrightarrow
 H^1\bigl(G,K_2\overline k(X)/H^0(\overline X,\mathcal K_2)\bigr)
 &\longrightarrow\Ker[\CH^2(X)\to\CH^2(\overline X)]\\
 &\longrightarrow H^1\bigl(G,H^1(\overline X,\mathcal K_2)\bigr).
\end{split}
\]
仅假设 $H^1(X,\mathcal O_X)=0$，左端群便指数有界；假设 $H^2(X,\mathcal O_X)=0$，右端群便指数有界。为证明这些结论，[CR1] 使用了 $K_2$ 的 Hilbert 第九十定理的上同调形式 [C1]，即
\[
 H^1\bigl(G,K_2\overline k(X)/K_2\overline k\bigr)=0.
\]
它只要求 $X(k)\ne\varnothing$，无须对凝聚上同调群作任何假设。这样，左端群的研究便归结为研究
\[
 \Ker\bigl[H^2(G,H^0(\overline X,\mathcal K_2))
       \longrightarrow H^2(G,K_2\overline k(X))\bigr].
\]
同时，再利用 Merkur'ev 和 Suslin 的若干基本结果，分析 $H^0(\overline X,\mathcal K_2)$ 与 $H^1(\overline X,\mathcal K_2)$ 的结构。在 $H^1(X,\mathcal O_X)=0$ 的假设下，可以证明 $H^2(G,H^0(\overline X,\mathcal K_2))$ 本身就指数有界。详情见 [CR1] 及综述 [R1]。\hfill$\square$

一方面使用更完整的\term{局部化正合列}{localization exact sequence}{suite exacte de localisation}，另一方面使用由 Kato--Saito 和 Somekawa 发展的 Bloch 的一项基本结果，便可去掉定理 6.3 中的条件 $H^1(X,\mathcal O_X)=0$。

\begin{statement}{定理 6.5}
设 $X$ 是\term{数域}{number field}{corps de nombres} $k$ 上的光滑、射影、几何连通簇。假设 $H^2(X,\mathcal O_X)=0$。设 $n>0$ 为自然数，则 $n$-挠子群 ${}_n\CH^2(X)$ 有限。
\end{statement}
\begin{proof}
 a) 存在有限域扩张 $L/k$，使 $X(L)\ne\varnothing$，且有限生成群之间的自然映射 $\NS(X_L)\to\NS(\overline X)$ 为同构。令 $m=[L:k]$。同定理 6.3 的证明一样，可以找到 $k$ 的整数环 $R$ 的一个局部化 $A=R_f$（$f\in R$，$f\ne0$），以及以 $X/k$ 为一般纤维的光滑射影 $A$-概形 $\mathcal X$，使 $m,n$ 在 $A$ 中可逆，$A$ 在 $L$ 中的整闭包 $B$ 在 $A$ 上有限平展，且理想类群 $\Pic(B)$ 为零（即 $B$ 为主理想环）。借助 Grothendieck 的半连续性定理，还可使每个闭点 $p\in\Spec(A)$ 上的纤维 $X_p$ 都满足 $H^2(X_p,\mathcal O_{X_p})=0$；于是，对 $\mathcal X_B/\Spec(B)$ 在闭点 $q\in\Spec(B)$ 上的纤维，同样的结论也成立。

\origpage{31}{26}
b) 定理 6.3 中使用的局部化正合列可以向左延伸为
\[
 H^1(X_L,\mathcal K_2)\longrightarrow
 \bigoplus_{q\in\Spec(B)^{(1)}}\Pic(X_q)
 \longrightarrow\CH^2(\mathcal X_B)
 \longrightarrow\CH^2(X_L)\longrightarrow0.
\]
事实上，它还可以继续向左延伸，见 [R2]，命题 1.2。\term{杯积}{cup product}{cup-produit}给出映射
\[
 \Pic(X_L)\otimes L^*\longrightarrow H^1(X_L,\mathcal K_2),
\]
利用层 $\mathcal K_2$ 的 Gersten--Quillen 分解，很容易描述这个映射。因此有复合映射 $\rho$：
\[
\begin{split}
 \Pic(\mathcal X_B)\otimes L^*
 &\longrightarrow\Pic(X_L)\otimes L^*
 \longrightarrow H^1(X_L,\mathcal K_2)\\
 &\longrightarrow\bigoplus_{q\in\Spec(B)^{(1)}}\Pic(X_q).
\end{split}
\]
容易验证，它把 $\mathcal L\otimes\alpha\in\Pic(\mathcal X_B)\otimes L^*$ 的类映为族 $(\mathcal L|_{X_q})^{v_q(\alpha)}$。同 6.3 中类似的论证，结合 Grothendieck 的提升定理，表明将 $\rho$ 与投影
\[
 \bigoplus_{q\in\Spec(B)^{(1)}}\Pic(X_q)
 \longrightarrow\bigoplus_{q\in\Spec(B)^{(1)}}\NS(X_q)
\]
复合，便得到满射。

因此复合映射
\[
 H^1(X_L,\mathcal K_2)\longrightarrow
 \bigoplus_{q\in\Spec(B)^{(1)}}\Pic(X_q)
 \longrightarrow\bigoplus_{q\in\Spec(B)^{(1)}}\NS(X_q)
\]
是满射；这里 $H^2(X,\mathcal O_X)=0$ 的假设起了关键作用。不过，要像 6.3 中那样完成证明，还须控制映射
\[
 H^1(X_L,\mathcal K_2)\longrightarrow
 \bigoplus_{q\in\Spec(B)^{(1)}}\Pic(X_q)
\]
的余核。也就是说，我们希望群
\[
 \bigoplus_{q\in\Spec(B)^{(1)}}\Pic^0(X_q)
\]
中的元素也能作为 $H^1(X_L,\mathcal K_2)$ 中元素的边界而出现。这就要用到下面回顾的 Bloch、Kato--Saito 和 Somekawa 的结果。

c) 设 $R$ 是 $p$-进域 $K$（即 $\mathbb Q_p$ 的有限扩张）的整数环，赋值为 $v:K^*\to\mathbb Z$，剩余域为 $\mathbb F$。设 $J$ 是 $R$ 上的阿贝尔概形。由 $J(R)=J(K)$，可定义\term{特殊化映射}{specialization map}{application de spécialisation} $J(K)\to J(\mathbb F)$，记为 $x\mapsto\overline x$。进而定义
\[
 J(K)\otimes K^*\longrightarrow J(\mathbb F),
 \qquad x\otimes\lambda\longmapsto v(\lambda)\,\overline x.
\]
若 $M/K$ 是有限域扩张，定义映射
\[
 \rho_{M/K}:J(M)\otimes M^*\longrightarrow J(\mathbb F)
\]
\origpage{32}{27}
为复合映射
\[
 J(M)\otimes M^*\longrightarrow J(\mathbb F_M)
 \xrightarrow{N}J(\mathbb F),
\]
其中 $N$ 是剩余域扩张 $\mathbb F_M/\mathbb F$ 的范数。

现在设 $K$ 是数域，$R=A_f$ 是其整数环 $A$ 的一个局部化（$f\in A$，$f\ne0$），$J$ 是 $R$ 上的阿贝尔概形。对每个闭点 $p\in\Spec(R)$，设其剩余域为 $\mathbb F_p$（这里 $p$ 是闭点，而不是素数！）。对每个有限扩张 $M/K$，将对应于 $p$ 在 $R$ 的 $M$ 中整闭包中的各个延拓 $q$ 的映射 $\rho_{M_q/K_p}$ 相加，便定义了
\[
 \rho_{M/K,p}:J(M)\otimes M^*\longrightarrow J(\mathbb F_p).
\]
我们要使用的关键结果如下。其思想源自 Bloch [B5]，后来由 Kato--Saito [KS]、Somekawa [So] 加以发展。

\begin{statement}{定理}
在上述情形中，映射
\[
 \bigoplus_{M/K}\bigoplus_p\rho_{M/K,p}:
 \bigoplus_{M/K}\bigl(J(M)\otimes M^*\bigr)
 \longrightarrow\bigoplus_p J(\mathbb F_p)
\]
的余核是有限群。这里 $M/K$ 遍历 $K$ 的有限扩张，$p$ 遍历 $R$ 的素理想。
\end{statement}

d) 利用 $H^2(X_q,\mathcal O_{X_q})=0$，并在必要时进一步缩小 $A$，由 Grothendieck 的另一个结果，可使\term{Picard 函子}{Picard functor}{foncteur de Picard} $\underline{\Pic}_{\mathcal X_B/B}$ 由一个局部光滑的 $B$-群概形表示。其单位连通分支 $J=\underline{\Pic}^{\,0}_{\mathcal X_B/B}$ 是一个 $B$-阿贝尔概形。

考虑下图：
\[
\begin{tikzcd}[ampersand replacement=\&,column sep=large,row sep=large]
 \& \displaystyle\bigoplus_{q\in\Spec(B)^{(1)}}\NS(X_q)\\
 H^1(X_L,\mathcal K_2) \arrow[r]\arrow[ur]
 \& \displaystyle\bigoplus_{q\in\Spec(B)^{(1)}}\Pic(X_q)\arrow[u]\\
 \displaystyle\bigoplus_{M/L}\bigl(J(M)\otimes M^*\bigr)\arrow[u]\arrow[r]
 \& \displaystyle\bigoplus_{q\in\Spec(B)^{(1)}}J(\mathbb F_q)\arrow[u]
\end{tikzcd}
\]
其定义如下。$M$ 遍历 $L$ 的有限扩张，每个同构类取一个代表。中间的水平箭头来自局部化正合列。左侧的竖直箭头由复合同态
\[
\begin{split}
 J(M)\otimes M^*&\longrightarrow\Pic(X_M)\otimes M^*
 \xrightarrow{\mathrm{cup}}H^1(X_M,\mathcal K_2)\\
 &\xrightarrow{N_{M/L}}H^1(X_L,\mathcal K_2)
\end{split}
\]
定义。斜箭头按定义使三角形交换。最下方的水平箭头就是上面 c) 中描述的映射。可以验证，该方块交换。

由 b)，斜箭头为满射；由 c)，最下方水平箭头的余核有限。因此中间水平箭头的余核也有限。

\origpage{33}{28}
e) 从 b) 的局部化正合列可知，映射
\[
 \CH^2(\mathcal X_B)\longrightarrow\CH^2(X_L)
\]
为满射，且核有限。再用定理 6.2，并仿照 6.3 证明末尾作图追踪，即可推出 ${}_n\CH^2(X)$ 有限。
\end{proof}

由定理 6.5 出发，要在完全一般的情形下证明定理 6.1，还须去掉定理 6.4 中的条件 $H^1(X,\mathcal O_X)=0$。换言之，需要证明：

\begin{statement}{定理 6.6（Salberger）}
设 $X$ 是数域 $k$ 上的光滑、射影、几何连通簇。假设 $H^2(X,\mathcal O_X)=0$，则挠群 $\CH^2(X)_{\tors}$ 指数有界。
\end{statement}
沿用定理 6.4 的方法可知，这可由下述命题推出。
\begin{statement}{定理 6.7（Salberger）}
设 $X$ 是数域 $k$ 上的光滑、射影、几何连通簇，则群
\[
 \Ker\rho_X:=\Ker\bigl[H^2(G,H^0(\overline X,\mathcal K_2))
 \longrightarrow H^2(G,K_2\overline k(X))\bigr]
\]
指数有界。
\end{statement}
\begin{proof}[证明概要]
Salberger 的关键思想是，逐次取光滑超平面截面，将命题归结为曲线情形；Bertini 定理保证了这种截面的存在。设 $Y\subset X$ 为这样的截面。可以证明，限制映射
\[
 H^2(G,H^0(\overline X,\mathcal K_2))
 \longrightarrow H^2(G,H^0(\overline Y,\mathcal K_2))
\]
诱导映射 $\Ker(\rho_X)\to\Ker(\rho_Y)$。另一方面，根据 [CT--R 1985] 中基于 Suslin 结果的论证，有正合列
\[
 0\longrightarrow H^1_{\et}(\overline X,\mathbb Q/\mathbb Z(2))
 \longrightarrow H^0(\overline X,\mathcal K_2)
 \longrightarrow H^0(\overline X,\mathcal K_2)\otimes\mathbb Q
 \longrightarrow0,
\]
从而有同构
\[
 H^2\bigl(G,H^1_{\et}(\overline X,\mathbb Q/\mathbb Z(2))\bigr)
 \simeq H^2\bigl(G,H^0(\overline X,\mathcal K_2)\bigr).
\]
将限制映射限制到 $G$-模 $H^1_{\et}(\overline X,\mathbb Q/\mathbb Z(2))$ 与 $H^1_{\et}(\overline Y,\mathbb Q/\mathbb Z(2))$ 的最大可除子群，得到
\[
 H^1_{\et}(\overline X,\mathbb Q/\mathbb Z(2))^0
 \longrightarrow H^1_{\et}(\overline Y,\mathbb Q/\mathbb Z(2))^0.
\]
这是一个 $G$-等变映射，且具有一个\term{准收缩}{almost retraction}{presque rétraction}：存在 $G$-等变映射
\[
 H^1_{\et}(\overline Y,\mathbb Q/\mathbb Z(2))^0
 \longrightarrow H^1_{\et}(\overline X,\mathbb Q/\mathbb Z(2))^0,
\]
它与前一映射复合后等于乘以某个正整数。在 [CR3] 中，我们利用 Poincaré 的完全可约性定理证明了这一点。$H^1_{\et}(\overline X,\mathbb Q/\mathbb Z(1))^0$ 这一最大可除子群正是 $X$ 的 Picard 簇的挠点群，因而可将上述定理应用于阿贝尔簇之间的单同态
\[
 \underline{\Pic}^{\,0}_{X/k}\longrightarrow\underline{\Pic}^{\,0}_{Y/k}.
\]
所以，只须处理曲线情形。对曲线，定理由 Raskind [R3] 证明；他使用了 Jannsen 的结果（数域上曲线类域论的局部整体原理 [J]）以及 S.~Saito 的结果（局部域上曲线的类域论）。事实上，在曲线情形，仅假设 $X(k)\ne\varnothing$，就有更强的结论：$\Ker(\rho_X)=0$。
\end{proof}

\origpage{34}{29}
\section*{§7. 数域上的簇，II}
\addcontentsline{toc}{section}{§7. 数域上的簇，II}
本节介绍 Shuji Saito [S4] 对上一节所述挠元有限性定理的处理方法。这一方法有时还能得到更进一步的结论（见下面的定理 7.6）。如前所述，Saito 的文章吸收了 Salberger 和我的一些建议。Saito 方法的一个有趣之处是：在一定假设下，它能保证取值于\term{平展上同调}{étale cohomology}{cohomologie étale}的\term{闭链映射}{cycle map}{application cycle}在 Chow 群的挠子群上为单射。

下面的叙述将突出 [S4] 中一些不太显眼的要点。

设 $X$ 是特征零域 $k$ 上的光滑整簇。对任意整数 $n>0$，有正合列 (3.11)：
\[
 0\longrightarrow H^1(X,\mathcal K_2)/n
 \longrightarrow NH^3_{\et}(X,\mu_n^{\otimes2})
 \longrightarrow {}_n\CH^2(X)\longrightarrow0,
\]
其中按定义，
\[
 NH^3_{\et}(X,\mu_n^{\otimes2})
 =\Ker[H^3_{\et}(X,\mu_n^{\otimes2})
       \longrightarrow H^3(k(X),\mu_n^{\otimes2})].
\]
对所有整数 $n>0$ 取归纳极限，得到正合列
\[
 0\longrightarrow H^1(X,\mathcal K_2)\otimes\mathbb Q/\mathbb Z
 \longrightarrow NH^3_{\et}(X,\mathbb Q/\mathbb Z(2))
 \longrightarrow\CH^2(X)_{\tors}\longrightarrow0,
\]
其中
\[
 NH^3_{\et}(X,\mathbb Q/\mathbb Z(2))
 =\Ker[H^3_{\et}(X,\mathbb Q/\mathbb Z(2))
       \longrightarrow H^3(k(X),\mathbb Q/\mathbb Z(2))].
\]
记
\[
 D=H^0(X,\mathcal H^3(\mathbb Q/\mathbb Z(2)))
 \subset H^3(k(X),\mathbb Q/\mathbb Z(2)),
\]
并记
\[
 \Theta=\coker[H^1(X,\mathcal K_2)\otimes\mathbb Q/\mathbb Z
                \longrightarrow H^3_{\et}(X,\mathbb Q/\mathbb Z(2))].
\]
于是得到下列显然交换的正合列图：
\begin{equation*}
\begin{tikzcd}[ampersand replacement=\&,column sep=small,row sep=large]
 \&\&0\arrow[d]\&0\arrow[d]\\
 0\arrow[r]\&H^1(X,\mathcal K_2)\otimes\mathbb Q/\mathbb Z\arrow[r]\arrow[d,equal]
 \&NH^3_{\et}(X,\mathbb Q/\mathbb Z(2))\arrow[r]\arrow[d]
 \&\CH^2(X)_{\tors}\arrow[r]\arrow[d]\&0\\
 0\arrow[r]\&H^1(X,\mathcal K_2)\otimes\mathbb Q/\mathbb Z\arrow[r]
 \&H^3_{\et}(X,\mathbb Q/\mathbb Z(2))\arrow[r]\arrow[d]
 \&\Theta\arrow[r]\arrow[dl,"{\Psi}"']\&0\\
 \&\&D
\end{tikzcd}
\tag{7.2}
\end{equation*}
这里 $\Psi$ 是由下方竖直箭头诱导的映射。

下述定理概括了 Saito 方法的代数部分。稍后将给出保证假设 (H) 成立的算术和几何条件。
\begin{statement}{定理 7.1}
设 $X$ 是特征零域 $k$ 上的光滑整簇。作如下假设：

\noindent (H) 群
\[
 \Theta=\coker[H^1(X,\mathcal K_2)\otimes\mathbb Q/\mathbb Z
        \longrightarrow H^3_{\et}(X,\mathbb Q/\mathbb Z(2))]
\]
指数有界，且被整数 $N>0$ 零化。
\origpage{35}{30}
则 $\CH^2(X)_{\tors}$ 的指数有限，且整除 $N$；对任意正整数 $n$，只要 $N\mid n$，自然投影与闭链映射的复合
\[
 \CH^2(X)_{\tors}\longrightarrow\CH^2(X)/n
 \longrightarrow H^4_{\et}(X,\mu_n^{\otimes2})
\]
就是单射。
\end{statement}
\begin{proof}
a) 第一项结论，即在假设 (H) 下 $N$ 消去 $\CH^2(X)_{\tors}$，可直接由图 (7.2) 看出：事实上，$\CH^2(X)_{\tors}$ 是 $\Theta$ 的子群。

b) 先证明一个引理。
\begin{statement}{引理（[S4]，命题 2.4）}
对任意光滑整 $k$-簇 $X$ 和任意整数 $n>0$，有正合列
\[
 {}_n\Theta\longrightarrow {}_nD\longrightarrow\CH^2(X)/n
 \longrightarrow H^4_{\et}(X,\mu_n^{\otimes2}).
\]
\end{statement}
\noindent\textit{引理的证明。} 考虑下列交换图：
\begin{equation*}
\resizebox{.98\linewidth}{!}{$
\begin{tikzcd}[ampersand replacement=\&,column sep=small,row sep=large]
 \&\&0\arrow[d]\\
 \&H^3_{\et}(X,\mu_n^{\otimes2})\arrow[r]\arrow[d]
 \&H^0(X,\mathcal H^3(\mu_n^{\otimes2}))\arrow[r]\arrow[d]
 \&\CH^2(X)/n\arrow[r]\&H^4_{\et}(X,\mu_n^{\otimes2})\\
 H^1(X,\mathcal K_2)\otimes\mathbb Q/\mathbb Z\arrow[r]
 \&H^3_{\et}(X,\mathbb Q/\mathbb Z(2))\arrow[r]\arrow[d,"{\times n}"']
 \&H^0(X,\mathcal H^3(\mathbb Q/\mathbb Z(2)))\arrow[d,"{\times n}"]\\
 \&H^3_{\et}(X,\mathbb Q/\mathbb Z(2))\arrow[r]
 \&H^0(X,\mathcal H^3(\mathbb Q/\mathbb Z(2))).
\end{tikzcd}$}
\tag{7.3}
\end{equation*}
图中，上方的水平列是来自 Bloch--Ogus 理论的正合列（见 §3）；中间的水平列是一个复形（见本节开头）；左侧的竖直列正合，因为它是由平展层正合列
\[
 1\longrightarrow\mu_n^{\otimes2}\longrightarrow\mathbb Q/\mathbb Z(2)
 \xrightarrow{\times n}\mathbb Q/\mathbb Z(2)\longrightarrow1
\]
得到的 Kummer 正合列。

下面证明右侧竖直列正合。考虑交换图
\[
\resizebox{.98\linewidth}{!}{$
\begin{tikzcd}[ampersand replacement=\&,column sep=large,row sep=large]
 \&0\arrow[d]\&0\arrow[d]\&0\arrow[d]\\
 0\arrow[r]\&H^0(X,\mathcal H^3(\mu_n^{\otimes2}))\arrow[r]\arrow[d]
 \&H^3_{\et}(k(X),\mu_n^{\otimes2})\arrow[r]\arrow[d]
 \&\displaystyle\bigoplus_{x\in X^1}H^2_{\et}(k(x),\mu_n)\arrow[d]\\
 0\arrow[r]\&H^0(X,\mathcal H^3(\mathbb Q/\mathbb Z(2)))\arrow[r]\arrow[d,"{\times n}"']
 \&H^3_{\et}(k(X),\mathbb Q/\mathbb Z(2))\arrow[r]\arrow[d,"{\times n}"]
 \&\displaystyle\bigoplus_{x\in X^1}H^2_{\et}(k(x),\mathbb Q/\mathbb Z(1))\arrow[d,"{\times n}"]\\
 0\arrow[r]\&H^0(X,\mathcal H^3(\mathbb Q/\mathbb Z(2)))\arrow[r]
 \&H^3_{\et}(k(X),\mathbb Q/\mathbb Z(2))\arrow[r]
 \&\displaystyle\bigoplus_{x\in X^1}H^2_{\et}(k(x),\mathbb Q/\mathbb Z(1)).
\end{tikzcd}$}
\]
\origpage{36}{31}
其中各水平列正合。竖直列来自 Kummer 正合列，所以除可能在首项外都正合。$H^2_{\et}(k(x),\mu_n)$ 到 $H^2_{\et}(k(x),\mathbb Q/\mathbb Z(1))$ 的映射为单射，是 Hilbert 第九十定理的推论。而 $H^3_{\et}(k(X),\mu_n^{\otimes2})$ 到 $H^3_{\et}(k(X),\mathbb Q/\mathbb Z(2))$ 的映射为单射，则是高度非平凡的结论；它由 Merkur'ev--Suslin 的主定理 [MS] 推出。

由上述图的正合性可得，图 (7.3) 右侧的竖直列正合。再由图 (7.3) 得到正合列
\[
 H^3_{\et}(X,\mu_n^{\otimes2})\longrightarrow{}_nD
 \longrightarrow\CH^2(X)/n\longrightarrow H^4_{\et}(X,\mu_n^{\otimes2}).
\]
由于首个箭头经由 ${}_n\Theta$ 分解（见 (7.2)），便得到所需的正合列
\[
 {}_n\Theta\longrightarrow{}_nD\longrightarrow\CH^2(X)/n
 \longrightarrow H^4_{\et}(X,\mu_n^{\otimes2}).
\]
引理得证。\hfill$\square$

c) 考虑下图：
\[
\begin{tikzcd}[ampersand replacement=\&,column sep=large,row sep=large]
 \&\&\CH^2(X)_{\tors}\arrow[d]\arrow[dr,bend left=15,"{\rho}"]\\
 {}_n\Theta\arrow[r]\arrow[d]\&{}_nD\arrow[r]\arrow[d]
 \&\CH^2(X)/n\arrow[r]\arrow[d]\&H^4_{\et}(X,\mu_n^{\otimes2})\\
 \Theta\arrow[r]\&D\arrow[r]\&\CH^2(X)\otimes\mathbb Q/\mathbb Z.
\end{tikzcd}
\]
其中，中间的水平列是刚才得到的正合列；下方水平列由这些正合列取归纳极限而得；$\rho$ 是复合映射。注意：
\begin{itemize}
\item 右侧通往 $\CH^2(X)\otimes\mathbb Q/\mathbb Z$ 的竖直列是复形。
\item 若 $N\mid n$，则 $n$ 消去 $\Theta$，因而 ${}_n\Theta\to\Theta$ 为满射。
\item 若 $N\mid n$，如前所述，$n$ 也消去 $\CH^2(X)_{\tors}$，因此映射 $\CH^2(X)_{\tors}\to\CH^2(X)/n$ 为单射。
\end{itemize}
固定满足 $N\mid n$ 的整数 $n>0$，设 $z_0\in\CH^2(X)_{\tors}$ 且 $\rho(z_0)=0$。通过图追踪即可得到 $z_0=0$，从而证明定理。下面给出这一追踪的示意图，元素 $z_i$ 的下标按其出现顺序排列：
\[
\begin{tikzcd}[ampersand replacement=\&,column sep=large,row sep=large]
 \&\&z_0\arrow[d,"{\rho}"]\\
 z_5\arrow[r]\arrow[d]\&z_6\;z_2\arrow[r]\arrow[d,shift left=2pt]\arrow[d,shift right=2pt]
 \&z_1\arrow[r]\arrow[d]\&0\\
 z_4\arrow[r]\&z_3\arrow[r]\&0.
\end{tikzcd}
\]
由单射 ${}_nD\subset D$ 可知 $z_6=z_2$，因此 $z_1=0$，最终得到 $z_0=0$。
\end{proof}

\origpage{37}{32}
\noindent\textbf{备注 7.1.1。} 上述定理中证明单射性的方法，首次出现在 P.~Salberger 一项尚在撰写中的工作里。在那里，Salberger 研究 $p$-进域上光滑射影曲面 $X$ 的群 $SK_1(X)$ 上的互反映射，并在一定假设下证明了一个单射性结论。Saito 则提出将同一技巧用于挠闭链的研究。

定理 7.1 有如下变体，适用于基域 $k$ 的\term{上同调维数}{cohomological dimension}{dimension cohomologique}大于 $2$ 的情形。
\begin{statement}{定理 7.2}
设 $X$ 是特征零域 $k$ 上的光滑整簇。作如下假设：

\noindent (H$'$) 群
\[
 \Theta'=\coker\left[
 H^3_{\et}(k,\mathbb Q/\mathbb Z(2))\oplus
 \bigl(H^1(X,\mathcal K_2)\otimes\mathbb Q/\mathbb Z\bigr)
 \longrightarrow H^3_{\et}(X,\mathbb Q/\mathbb Z(2))\right]
\]
指数有界，且被整数 $N>0$ 零化。

再假设 $X(k)\ne\varnothing$。则 $\CH^2(X)_{\tors}$ 的指数有限，且整除 $N$；对任意满足 $N\mid n$ 的整数 $n>0$，自然投影与闭链映射的复合
\[
 \CH^2(X)_{\tors}\longrightarrow\CH^2(X)/n
 \longrightarrow H^4_{\et}(X,\mu_n^{\otimes2})
\]
为单射。
\end{statement}
\begin{proof}[证明概要]
选取一点 $P\in X(k)$。该点给出了定义 $\Theta'$ 时使用的基变换映射
\[
 H^3_{\et}(k,\mathbb Q/\mathbb Z(2))
 \longrightarrow H^3_{\et}(X,\mathbb Q/\mathbb Z(2))
\]
的一个截面，从而可写成 $\Theta=\Theta'\oplus H^3_{\et}(k,\mathbb Q/\mathbb Z(2))$。利用同一点 $P$，还可将 $D=H^0(X,\mathcal H^3(\mathbb Q/\mathbb Z(2)))$ 分解为
\[
 D=D'\oplus H^3_{\et}(k,\mathbb Q/\mathbb Z(2)).
\]
于是得到正合列
\[
 {}_n\Theta'\longrightarrow{}_nD'\longrightarrow\CH^2(X)/n
 \longrightarrow H^4_{\et}(X,\mu_n^{\otimes2}),
\]
其余证明与定理 7.1 相同。
\end{proof}

当 $X$ 是特征零域 $k$ 上的光滑、射影、几何整簇时，有（[CR1]）
\[
 H^1(\overline X,\mathcal K_2)\otimes\mathbb Q/\mathbb Z=0.
\]
因此，余核为 $\Theta'$ 的映射
\[
 H^3_{\et}(k,\mathbb Q/\mathbb Z(2))\oplus
 \bigl(H^1(X,\mathcal K_2)\otimes\mathbb Q/\mathbb Z\bigr)
 \longrightarrow H^3_{\et}(X,\mathbb Q/\mathbb Z(2))
\]
可经由下列映射分解：
\[
\begin{split}
 H^3_{\et}(k,\mathbb Q/\mathbb Z(2))\oplus
 \bigl(H^1(X,\mathcal K_2)\otimes\mathbb Q/\mathbb Z\bigr)
 \longrightarrow\Ker[H^3_{\et}(X,\mathbb Q/\mathbb Z(2))
 \to H^3_{\et}(\overline X,\mathbb Q/\mathbb Z(2))].
\end{split}
\]
从 Salberger 的报告（见 §9）和 Saito 的文章，还可提取下述结果；它不要求对 $k$ 作任何算术假设。
\begin{statement}{定理 7.3}
设 $X$ 是特征零域 $k$ 上的光滑、射影、几何整簇。假设 $H^1(X,\mathcal O_X)=0$ 且 $H^2(X,\mathcal O_X)=0$。则映射
\[
\begin{split}
 H^3_{\et}(k,\mathbb Q/\mathbb Z(2))\oplus
 \bigl(H^1(X,\mathcal K_2)\otimes\mathbb Q/\mathbb Z\bigr)
 \longrightarrow\Ker[H^3_{\et}(X,\mathbb Q/\mathbb Z(2))
 \to H^3_{\et}(\overline X,\mathbb Q/\mathbb Z(2))]
\end{split}
\]
的余核指数有界。
\end{statement}

\origpage{38}{33}
\begin{proof}
平展上同调的 Hochschild--Serre 谱序列
\[
 H^p\bigl(k,H^q_{\et}(\overline X,\mathbb Q/\mathbb Z(2))\bigr)
 \Longrightarrow H^{p+q}_{\et}(X,\mathbb Q/\mathbb Z(2))
\]
在群
\[
 F^2=\Ker[H^3_{\et}(X,\mathbb Q/\mathbb Z(2))
 \longrightarrow H^3_{\et}(\overline X,\mathbb Q/\mathbb Z(2))]
\]
上给出过滤 $F^0\subset F^1\subset F^2$。有满射 $H^3_{\et}(k,\mathbb Q/\mathbb Z(2))\to F^0$，它与包含映射 $F^0\subset F^2$ 的复合正是自然映射 $H^3_{\et}(k,\mathbb Q/\mathbb Z(2))\to H^3_{\et}(X,\mathbb Q/\mathbb Z(2))$。商群 $F^1/F^0$ 是
\[
 H^2\bigl(k,H^1_{\et}(\overline X,\mathbb Q/\mathbb Z(2))\bigr)
\]
的子商群。在 $H^1(X,\mathcal O_X)=0$ 的假设下，$H^1_{\et}(\overline X,\mathbb Q/\mathbb Z(2))$ 为有限群，因此 $F^1/F^0$ 指数有界。商群 $F^2/F^1$ 是 $H^1(k,H^2_{\et}(\overline X,\mathbb Q/\mathbb Z(2)))$ 的子群。因此，只须证明：一旦 $H^2(X,\mathcal O_X)=0$，群
\[
 \coker\left[H^1(X,\mathcal K_2)\otimes\mathbb Q/\mathbb Z
 \longrightarrow H^1\bigl(k,H^2_{\et}(\overline X,\mathbb Q/\mathbb Z(2))\bigr)\right]
\]
就指数有界。

我们只要求指数有限，因而可用转移论证把 $k$ 换为一个有限扩张，仍记为 $k$，并假设 $X(k)\ne\varnothing$，且伽罗瓦群 $G=\Gal(\overline k/k)$ 在 Néron--Severi 群 $\NS(\overline X)$ 上的作用平凡（也可以在证明开头就作这一步）。于是 $\NS(X)=\NS(\overline X)$。考虑下图，其中 $t$ 表示相应群的挠子群：
\[
\resizebox{.98\linewidth}{!}{$
\begin{tikzcd}[ampersand replacement=\&,column sep=large,row sep=large]
 \Pic(X)\otimes k^*\otimes\mathbb Q/\mathbb Z\arrow[r]\arrow[d,equal]
 \&H^1(X,\mathcal K_2)\otimes\mathbb Q/\mathbb Z\arrow[r]
 \&H^1(k,H^2_{\et}(\overline X,\mathbb Q/\mathbb Z(2)))\\
 \Pic(X)\otimes H^1(k,\mathbb Q/\mathbb Z(1))\arrow[d]\&\&\\
 (\NS(X)/t)\otimes H^1(k,\mathbb Q/\mathbb Z(1))\arrow[r,"{\simeq}"]
 \&H^0(k,(\NS(\overline X)/t)\otimes H^1(k,\mathbb Q/\mathbb Z(1)))\arrow[r]
 \&H^1(k,(\NS(\overline X)/t)\otimes\mathbb Q/\mathbb Z(1))\arrow[uu,"{\varphi}"']
\end{tikzcd}$}
\]
首个水平箭头由杯积给出；$\varphi$ 由映射
\[
 \NS(\overline X)\otimes\mathbb Q/\mathbb Z(1)
 \longrightarrow H^2_{\et}(\overline X,\mathbb Q/\mathbb Z(2))
\]
诱导，而此映射又由 Kummer 正合列所得的映射
\[
 \delta:\NS(\overline X)\otimes\mathbb Q/\mathbb Z
 =\Pic(\overline X)\otimes\mathbb Q/\mathbb Z
 \longrightarrow H^2_{\et}(\overline X,\mathbb Q/\mathbb Z(1))
\]
作 $\mathbb Q/\mathbb Z(1)$ 扭转而来。根据 Saito，上图交换。由于处于特征零情形，Hodge 理论说明，条件 $H^2(X,\mathcal O_X)=0$ 等价于 $\delta$ 的余核有限。于是 $\coker(\varphi)$ 指数有界，结合上图即得定理。
\end{proof}

\noindent\textbf{备注 7.3.1。} 证明中实际上只用到了 $\Pic(X)\otimes k^*$ 在 $H^1(X,\mathcal K_2)$ 中的像（扩张基域以后）。正如 Salberger 所指出的，若能在 $H^1(X,\mathcal K_2)$ 中构造出更有趣的元素，将十分有价值。

\noindent\textbf{备注 7.3.2。} 由等式
\[
 H^1(\overline X,\mathcal K_2)\otimes\mathbb Q/\mathbb Z=0,
\]
图 (7.2) 以及定理 7.3 可知：若 $X$ 是特征零域 $k$ 上的光滑、射影、几何整簇，且 $H^1(X,\mathcal O_X)=0$、
\origpage{39}{34}
$H^2(X,\mathcal O_X)=0$，则群 $\Ker[\CH^2(X)\to\CH^2(\overline X)]$ 指数有界。这正是定理 6.4 的证明中回顾的 [CR1] 的结果。

定理 7.3 有如下直接的算术应用。
\begin{statement}{推论 7.4}
设 $k$ 是有理数域 $\mathbb Q$ 上的有限生成域，$X$ 是 $k$ 上的光滑、射影、几何整簇。假设 $H^1(X,\mathcal O_X)=0$ 且 $H^2(X,\mathcal O_X)=0$。则映射
\[
 H^3_{\et}(k,\mathbb Q/\mathbb Z(2))\oplus
 \bigl(H^1(X,\mathcal K_2)\otimes\mathbb Q/\mathbb Z\bigr)
 \longrightarrow H^3_{\et}(X,\mathbb Q/\mathbb Z(2))
\]
的余核指数有界。
\end{statement}
\begin{proof}
事实上，映射
\[
 H^3_{\et}(X,\mathbb Q/\mathbb Z(2))
 \longrightarrow H^3_{\et}(\overline X,\mathbb Q/\mathbb Z(2))
\]
的像落在伽罗瓦群 $G=\Gal(\overline k/k)$ 的不变子群 $H^3_{\et}(\overline X,\mathbb Q/\mathbb Z(2))^G$ 中。正如定理 6.4 的证明中已经提到的，$k$ 在素域上有限生成这一假设，加上向有限域情形的约化，通过 Deligne 对 Weil 猜想的证明，保证这个群是有限的；更一般地，当 $i\ne2j$ 时，所有 $H^i_{\et}(\overline X,\mathbb Q/\mathbb Z(j))^G$ 都有限（[CR1]，定理 1.5）。由定理 7.3 立即得到结论。
\end{proof}

推论 7.4 说明，在其假设下，定理 7.2 的条件 (H$'$) 成立。

另一方面，如 1.9.2 中所预告的，有下述一般定理。
\begin{statement}{定理 7.5（[S4]）}
设 $X$ 是 $\mathbb Q$ 上有限生成域 $k$ 上的光滑簇，不要求其为固有簇。则对任意整数 $n>0$ 及任意自然数 $i$，闭链映射
\[
 \CH^i(X)/n\longrightarrow H^{2i}_{\et}(X,\mu_n^{\otimes i})
\]
的像是有限群。
\end{statement}
\begin{proof}[证明概要]
可以找到一个分式域为 $k$、在 $\mathbb Z$ 上光滑且有限型的环 $R$，使 $n$ 在 $R$ 中可逆；还可以找到一个在 $\Spec(R)$ 上光滑的概形 $\mathcal X$，使 $\mathcal X\times_Rk=X$。利用 $p:\mathcal X\to\Spec(\mathbb Z)$ 的光滑性，可知平展层 $R^jp_*\mu_n^{\otimes i}$ 在 $\Spec(\mathbb Z)$ 上可构造，故其平展上同调有限。再通过 Leray 谱序列，便知 $H^r_{\et}(\mathcal X,\mu_n^{\otimes s})$ 有限（见 6.2）。映射 $\CH^i(\mathcal X)\to\CH^i(X)$ 是满射，但目前尚不能定义从 $\CH^i(\mathcal X)$ 到 $H^{2i}_{\et}(\mathcal X,\mu_n^{\otimes i})$ 的闭链映射；若能定义并使自然的图交换，所需的有限性就会随之成立。Saito 证明中精细的部分在于说明：至少闭链映射 $\CH^i(X)\to H^{2i}_{\et}(X,\mu_n^{\otimes i})$ 的像，包含在 $H^{2i}_{\et}(\mathcal X,\mu_n^{\otimes i})$ 在 $H^{2i}_{\et}(X,\mu_n^{\otimes i})$ 中的像内，从而保证其有限性。
\end{proof}

结合 7.2、7.4 和 7.5，便得到下述有限性定理。
\begin{statement}{定理 7.6}
设 $k$ 是有理数域 $\mathbb Q$ 上的有限生成域，$X$ 是 $k$ 上的光滑、射影、几何整簇。假设 $H^1(X,\mathcal O_X)=0$、$H^2(X,\mathcal O_X)=0$，且 $X(k)\ne\varnothing$。则 $\CH^2(X)$ 的挠子群 $\CH^2(X)_{\tors}$ 有限。
\end{statement}

\origpage{40}{35}
\noindent\textbf{备注 7.6.1。} 除了额外要求 $X(k)\ne\varnothing$ 外，这一定理推广了 6.3.1 在数域上得到的结论；6.3.1 中并未施加这一条件。在 $\mathbb Q$ 上有限生成的域 $k$ 上，此前只处理过有理曲面情形（[C1]），且同样要求 $X(k)\ne\varnothing$。事实上，在 $p$-进域 $\mathbb Q_p$ 上的一元有理函数域 $k=\mathbb Q_p(T)$ 上，存在有理曲面 $X$，满足 $X(k)=\varnothing$ 且 $\CH^2(X)_{\tors}\ne0$，而闭链映射无法探测到这部分挠元（Sansuc 与作者，\textit{J.~of Algebra} \textbf{84} (1985)；Salberger 曾作出未发表的重新解释）。

Saito 运用自己的方法，一方面借助 Salberger 对挠元指数的有界性结果（定理 6.6），另一方面借助与定理 6.7 证明中所用结果相近的 Jannsen 的结果，成功得到了上一节中最一般的有限性定理。
\begin{statement}{定理 7.7（即定理 6.1）}
设 $X$ 是数域 $k$ 上的光滑、射影、几何连通簇。假设 $H^2(X,\mathcal O_X)=0$，则挠子群 $\CH^2(X)_{\tors}$ 有限。
\end{statement}
\begin{proof}
根据 6.6，群 $\CH^2(X)_{\tors}$ 指数有界。因此，对几乎所有素数 $l$，其 $l$-准素部分 $\CH^2(X)_{l\text{-}\tors}$ 为零；只须对每个素数 $l$ 证明 $\CH^2(X)_{l\text{-}\tors}$ 有限即可。如定理 7.3 中一样，考虑平展上同调谱序列
\[
 H^p\bigl(k,H^q_{\et}(\overline X,\mathbb Q_l/\mathbb Z_l(2))\bigr)
 \Longrightarrow H^{p+q}_{\et}(X,\mathbb Q_l/\mathbb Z_l(2)),
\]
它在 $H^3_{\et}(X,\mathbb Q_l/\mathbb Z_l(2))$ 上给出过滤
$F^0\subset F^1\subset F^2\subset F^3$。

在阿贝尔群范畴对有限阿贝尔群所取的\term{Serre 商范畴}{Serre quotient category}{catégorie quotient de Serre} $\mathcal C$ 中工作会十分方便。

由于 $k$ 为数域，$F^0$ 作为 $H^3(k,\mathbb Q_l/\mathbb Z_l(2))$ 的商群是有限的，所以在 $\mathcal C$ 中为零。$F^3/F^2$ 是有限群 $H^3_{\et}(\overline X,\mathbb Q_l/\mathbb Z_l(2))^G$ 的子群（见 7.4 的证明），所以在 $\mathcal C$ 中也为零。有单射
\[
 F^2/F^1\hookrightarrow H^1\bigl(k,H^2_{\et}(\overline X,\mathbb Q_l/\mathbb Z_l(2))\bigr).
\]
在 $\mathcal C$ 中还有满射
\[
 H^2\bigl(k,H^1_{\et}(\overline X,\mathbb Q_l/\mathbb Z_l(2))\bigr)
 \twoheadrightarrow F^1/F^0.
\]
在 $k$ 是数域的假设下，Jannsen（[J]，第 355 页，推论 7(b)）证明：存在整数 $r\ge0$（先验地依赖于 $l$，且对几乎所有 $l$ 可取 $r=0$），以及满射
\[
 (\mathbb Q_l/\mathbb Z_l)^r\twoheadrightarrow
 H^2\bigl(k,H^1_{\et}(\overline X,\mathbb Q_l/\mathbb Z_l(2))^0\bigr).
\]
这里 $M^0$ 表示群 $M$ 的最大可除子群。注意，商群
\[
 H^1_{\et}(\overline X,\mathbb Q_l/\mathbb Z_l(2))
 /H^1_{\et}(\overline X,\mathbb Q_l/\mathbb Z_l(2))^0
\]
有限，因此映射
\[
 H^2\bigl(k,H^1_{\et}(\overline X,\mathbb Q_l/\mathbb Z_l(2))^0\bigr)
 \longrightarrow H^2\bigl(k,H^1_{\et}(\overline X,\mathbb Q_l/\mathbb Z_l(2))\bigr)
\]
的余核指数有界。

有
\[
 F^2=\Ker[H^3_{\et}(X,\mathbb Q_l/\mathbb Z_l(2))
 \longrightarrow H^3_{\et}(\overline X,\mathbb Q_l/\mathbb Z_l(2))].
\]
如定理 7.3 的证明所示，在特征零域 $k$ 上，$H^2(X,\mathcal O_X)=0$ 蕴含复合映射
\[
 H^1(X,\mathcal K_2)\otimes\mathbb Q_l/\mathbb Z_l
 \longrightarrow F^2
 \longrightarrow H^1\bigl(k,H^2_{\et}(\overline X,\mathbb Q_l/\mathbb Z_l(2))\bigr)
\]
\origpage{41}{36}
的余核指数有界。

于是，在 $\mathcal C$ 中有复合映射
\[
\begin{split}
 (\mathbb Q_l/\mathbb Z_l)^r
 &\longrightarrow H^2\bigl(k,H^1_{\et}(\overline X,\mathbb Q_l/\mathbb Z_l(2))^0\bigr)\\
 &\longrightarrow H^2\bigl(k,H^1_{\et}(\overline X,\mathbb Q_l/\mathbb Z_l(2))\bigr)
 \longrightarrow H^3_{\et}(X,\mathbb Q_l/\mathbb Z_l(2)).
\end{split}
\]
此处记
\[
 \Theta=\coker[H^1(X,\mathcal K_2)\otimes\mathbb Q_l/\mathbb Z_l
 \longrightarrow H^3_{\et}(X,\mathbb Q_l/\mathbb Z_l(2))].
\]
结合图
\[
\resizebox{.98\linewidth}{!}{$
\begin{tikzcd}[ampersand replacement=\&,row sep=large,column sep=large]
 \&0\\
 \&\Theta\arrow[u]\\
 (\mathbb Q_l/\mathbb Z_l)^r\arrow[r]
 \&H^3_{\et}(X,\mathbb Q_l/\mathbb Z_l(2))\arrow[u]\arrow[r]
 \&H^1(k,H^2_{\et}(\overline X,\mathbb Q_l/\mathbb Z_l(2)))\\
 \&H^1(X,\mathcal K_2)\otimes\mathbb Q_l/\mathbb Z_l\arrow[u]
\end{tikzcd}$}
\]
及前述讨论可知，在当前假设下，群
\[
 \Theta''=\coker[(\mathbb Q_l/\mathbb Z_l)^r\longrightarrow\Theta]
\]
指数有界。再记（见图 (7.2)）
\[
 D''=\coker[(\mathbb Q_l/\mathbb Z_l)^r
 \longrightarrow H^0(X,\mathcal H^3(\mathbb Q_l/\mathbb Z_l(2)))].
\]
群 $(\mathbb Q_l/\mathbb Z_l)^a$ 的任意商群都形如 $(\mathbb Q_l/\mathbb Z_l)^b$（$b\le a$）。因此，在范畴 $\mathcal C$ 中，对每个整数 $n>0$，在 $n$-挠子群上诱导的自然映射 ${}_n\Theta\to{}_n\Theta''$ 和 ${}_nD\to{}_nD''$ 都是同构；并且存在整数 $n>0$，为 $l$ 的幂，使包含映射 ${}_n\Theta''\to\Theta''$ 为同构。同定理 7.1 证明中的 c) 一样，在 $\mathcal C$ 中考虑交换图
\[
\begin{tikzcd}[ampersand replacement=\&,column sep=large,row sep=large]
 \&\&\CH^2(X)_{l\text{-}\tors}\arrow[d]\\
 {}_n\Theta''\arrow[r]\arrow[d]\&{}_nD''\arrow[r]\arrow[d]
 \&\CH^2(X)/n\arrow[r]\arrow[d]\&H^4_{\et}(X,\mu_n^{\otimes2})\\
 \Theta''\arrow[r]\&D''\arrow[r]\&\CH^2(X)\otimes\mathbb Q_l/\mathbb Z_l.
\end{tikzcd}
\]
其中水平列正合（由 Merkur'ev--Suslin 定理），右侧竖直列是复形，且 ${}_n\Theta''\to\Theta''$ 为同构。根据 Salberger 的定理 6.6，$\CH^2(X)(l)$ 指数有界。因此，必要时将 $n$ 换为 $l$ 的更高次幂，使之消去 $\CH^2(X)_{l\text{-}\tors}$，便可保证映射 $\CH^2(X)_{l\text{-}\tors}\to\CH^2(X)/n$ 为单射。

与定理 7.1 相同的图追踪表明，映射
\[
 \CH^2(X)_{l\text{-}\tors}\longrightarrow H^4_{\et}(X,\mu_n^{\otimes2})
\]
\origpage{42}{37}
在 $\mathcal C$ 中为单射。因此，在阿贝尔群范畴中，映射
\[
 \CH^2(X)_{l\text{-}\tors}\longrightarrow H^4_{\et}(X,\mu_n^{\otimes2})
\]
的核有限。由于其像由定理 7.5 可知是有限的，遂得 $\CH^2(X)_{l\text{-}\tors}$ 本身有限。
\end{proof}

\noindent\textbf{备注 7.7.1。} 上述方法直接证明了 $\CH^2(X)_{\tors}$ 的有限性，无须像 §6 中那样，先对每个整数 $n>0$ 证明 ${}_n\CH^2(X)$ 有限。

\origpage{43}{38}
\section*{§8. 局部域上的簇}
\addcontentsline{toc}{section}{§8. 局部域上的簇}
这里的\term{局部域}{local field}{corps local}指 $p$-进域 $\mathbb Q_p$ 的有限扩张。有下述一般结果（参见 [CSS]，第 773 页，推论 2）。
\begin{statement}{定理 8.1}
设 $k$ 为局部域，$X$ 为光滑 $k$-簇。则对任意整数 $n>0$，$n$-挠子群 ${}_n\CH^2(X)$ 有限。
\end{statement}
\begin{proof}
这是定理 3.3.2 的直接推论：群 ${}_n\CH^2(X)$ 是平展上同调群 $H^3_{\et}(X,\mu_n^{\otimes2})$ 的子商群。而对局部域上的簇 $X$、任意整数 $i\ge0$ 和 $j$，群 $H^i_{\et}(X,\mu_n^{\otimes j})$ 都有限。这可由 Hochschild--Serre 谱序列、群 $H^q_{\et}(\overline X,\mu_n^{\otimes j})$（$q\ge0$）的有限性，以及当 $k$ 为局部域、$F$ 为有限 $\Gal(\overline k/k)$-模时群 $H^p(\Gal(\overline k/k),F)$（$p\ge0$）的有限性得到。
\end{proof}

下述定理在 [CR1] 中已于更强的假设下得到。下面的命题 8.3 使我们能够进一步去掉 $H^1(X,\mathcal O_X)=0$，或 $X$（或 $\Alb_X$）具有好约化之类的假设。
\begin{statement}{定理 8.2}
设 $k$ 为局部域，即 $\mathbb Q_p$ 的有限扩张，$X$ 为光滑、射影、几何整 $k$-簇。假设 $H^2(X,\mathcal O_X)=0$。则：
\begin{enumerate}[label=(\alph*)]
\item 群 $\Ker[\CH^2(X)\to\CH^2(\overline X)]$ 有限。
\item 若 $NH^3_{\et}(X,\mathbb Q/\mathbb Z(2))$ 在 $H^3_{\et}(\overline X,\mathbb Q/\mathbb Z(2))$ 中的像有限，例如 $H^3_{\et}(\overline X,\mathbb Q/\mathbb Z(2))^G$ 有限，则 $\CH^2(X)_{\tors}$ 有限。
\item （[CR1]，3.16）若 $X$ 在 $k$ 上具有\term{潜在好约化}{potentially good reduction}{potentielle bonne réduction}，则 $\CH^2(X)$ 中与 $p$ 互素的挠元构成有限群。
\item （[CR3]，6.1）若 $X$ 为曲面，则 $\CH^2(X)_{\tors}$ 有限。
\end{enumerate}
\end{statement}
\begin{proof}
首先，由定理 4.2，群 $\CH^2(\overline X)_{\tors}$ 以 $G$-等变方式嵌入 $H^3_{\et}(\overline X,\mathbb Q/\mathbb Z(2))$。观察图 (3.6) 可知，$\CH^2(X)_{\tors}$ 在
\[
 \CH^2(\overline X)_{\tors}\subset H^3_{\et}(\overline X,\mathbb Q/\mathbb Z(2))
\]
中的像，等同于
\[
 NH^3_{\et}(X,\mathbb Q/\mathbb Z(2))
 =\Ker[H^3_{\et}(X,\mathbb Q/\mathbb Z(2))
 \longrightarrow H^3_{\et}(k(X),\mathbb Q/\mathbb Z(2))]
\]
在 $H^3_{\et}(\overline X,\mathbb Q/\mathbb Z(2))$ 中的像。因此 (b) 由 (a) 推出。

若 $X$ 具有潜在好约化，由平展上同调的基变换定理和 Deligne 的定理（Weil 猜想），对每个素数 $l\ne p$，群 $H^3_{\et}(\overline X,\mathbb Q_l/\mathbb Z_l(2))^G$ 有限，且对几乎所有 $l$ 为零（[CR1]，1.5.1）。因此 (c) 由 (b) 推出。

若 $X$ 为曲面，则 $G$-模 $H^3_{\et}(\overline X,\mathbb Q/\mathbb Z(2))$ 等同于 $X$ 的\term{Albanese 簇}{Albanese variety}{variété d'Albanese} $\Alb_X$ 的挠点群 $\Alb_X(\overline k)_{\tors}$。于是 $H^3_{\et}(\overline X,\mathbb Q/\mathbb Z(2))^G$ 等同于 $\Alb_X(k)_{\tors}$；后者有限（见 §1，命题 1.7）。故 (d) 也由 (b) 推出。

下面证明 (a)。如定理 6.4 的证明中所述，有正合列
\[
\begin{split}
 H^1\bigl(G,K_2\overline k(X)/H^0(\overline X,\mathcal K_2)\bigr)
 &\longrightarrow\Ker[\CH^2(X)\to\CH^2(\overline X)]\\
 &\longrightarrow H^1\bigl(G,H^1(\overline X,\mathcal K_2)\bigr).
\end{split}
\]
在 $H^2(X,\mathcal O_X)=0$ 且 $k$ 为局部域的假设下，已知群 $H^1(G,H^1(\overline X,\mathcal K_2))$ 有限（[CR1]，命题 3.9）。利用局部域上簇的
\[
 H^1(G,K_2\overline k(X)/K_2\overline k)=0
\]
\origpage{44}{39}
（[C1]），以及群 $K_2\overline k$ 的唯一可除性，可知左端群等同于核
\[
 \Ker\rho_X:=\Ker[H^2(G,H^0(\overline X,\mathcal K_2))
 \longrightarrow H^2(G,K_2\overline k(X))].
\]
而下面的定理 8.3 说明该群有限。因而上述正合列的中间群也有限。
\end{proof}

\noindent\textbf{备注 8.2.1。} 可以问，(b) 的假设是否总能满足（参见 [J]，第 349 页，备注 5）。

因此，还须证明下述定理（[CR3]，命题 4.2）。它是 Salberger 的定理 6.7 在局部域上的对应结果。
\begin{statement}{定理 8.3}
设 $X$ 是局部域 $k$ 上的光滑、射影、几何连通簇。则群
\[
\begin{split}
 H^1\bigl(G,K_2\overline k(X)/H^0(\overline X,\mathcal K_2)\bigr)
 &=\Ker\rho_X\\
 &:=\Ker[H^2(G,H^0(\overline X,\mathcal K_2))
 \longrightarrow H^2(G,K_2\overline k(X))]
\end{split}
\]
是有限群。
\end{statement}
\begin{proof}
证明与定理 6.7 完全类似，沿用那里的记号。假设 $\dim X>1$，取光滑超平面截面 $Y\subset X$。限制映射
\[
 H^2(G,H^0(\overline X,\mathcal K_2))
 \longrightarrow H^2(G,H^0(\overline Y,\mathcal K_2))
\]
诱导映射 $\Ker(\rho_X)\to\Ker(\rho_Y)$。上述限制映射的核等同于
\[
 H^2(G,H^1_{\et}(\overline X,\mathbb Q/\mathbb Z(2)))
 \longrightarrow H^2(G,H^1_{\et}(\overline Y,\mathbb Q/\mathbb Z(2)))
\]
的核。由于有限 $G$-模的 $G$-上同调是有限的（$k$ 为局部域），忽略有限群后，该核又等同于
\[
 H^2(G,H^1_{\et}(\overline X,\mathbb Q/\mathbb Z(2))^0)
 \longrightarrow H^2(G,H^1_{\et}(\overline Y,\mathbb Q/\mathbb Z(2))^0)
\]
的核。

限制映射
\[
 H^1_{\et}(\overline X,\mathbb Q/\mathbb Z(2))^0
 \longrightarrow H^1_{\et}(\overline Y,\mathbb Q/\mathbb Z(2))^0
\]
为 $G$-等变映射，且有准收缩。也就是说，存在 $G$-等变映射
\[
 H^1_{\et}(\overline Y,\mathbb Q/\mathbb Z(2))^0
 \longrightarrow H^1_{\et}(\overline X,\mathbb Q/\mathbb Z(2))^0,
\]
它与前一映射复合后等于乘以某个整数 $m>0$。因此
\[
 H^2(G,H^1_{\et}(\overline X,\mathbb Q/\mathbb Z(2))^0)
 \longrightarrow H^2(G,H^1_{\et}(\overline Y,\mathbb Q/\mathbb Z(2))^0)
\]
的核包含在 $H^2(G,H^1_{\et}(\overline X,\mathbb Q/\mathbb Z(2))^0)$ 的 $m$-挠子群中，而该子群是
\[
 H^2\bigl(G,{}_m(H^1_{\et}(\overline X,\mathbb Q/\mathbb Z(2))^0)\bigr)
\]
的商群；由于 $k$ 为局部域，后者有限。所以 $\Ker(\rho_X)\to\Ker(\rho_Y)$ 的核有限。这样便归结为局部域上曲线 $C$ 的情形。Raskind 已在 [R3] 中完成了这一证明，下面简述其论证。曲线 $C$ 满足 $\CH^2(C)=0$，故伽罗瓦方法表明，群
\[
 H^1(G,K_2\overline k(C)/H^0(\overline C,\mathcal K_2))
\]
等同于映射
\[
 H^1(C,\mathcal K_2)\longrightarrow H^1(\overline C,\mathcal K_2)^G
\]
的余核。

\origpage{45}{40}
记
\[
 V(C)=\Ker[H^1(C,\mathcal K_2)\to k^*],\qquad
 V(\overline C)=\Ker[H^1(\overline C,\mathcal K_2)\to\overline k^{\,*}],
\]
其中箭头来自通常的互反映射。如 Raskind 所解释的（同上文献，3.5 及随后引理），S.~Saito 在局部域上的高维类域论中的一个定理保证了 $V(C)\to V(\overline C)^G$ 为满射。利用 $V(\overline C)$ 的唯一可除性（[R3]，引理 1.1），并对正合列
\[
 0\longrightarrow V(\overline C)\longrightarrow H^1(\overline C,\mathcal K_2)
 \longrightarrow\overline k^{\,*}\longrightarrow0
\]
取 $G$-上同调，可知 $H^1(C,\mathcal K_2)\to H^1(\overline C,\mathcal K_2)^G$ 的余核等同于
\[
 H^1(C,\mathcal K_2)\longrightarrow k^*
\]
的余核，也就是映射
\[
 N:\bigoplus_{P\in C^{(0)}}k(P)^*\longrightarrow k^*
\]
的余核。这里每个闭点 $P$ 对应的映射 $N_P:k(P)^*\to k^*$ 就是范数映射。若 $C$ 有 $k$-有理点，则 $N$ 为满射，此时甚至有
\[
 H^1(G,K_2\overline k(C)/H^0(\overline C,\mathcal K_2))=0.
\]
没有有理点时，$N$ 的像在 $k^*$ 中仍显然具有有限指数，因为 $k$ 为局部域，每个 $N_P$ 的像都已具有有限指数。可以问，即使没有 $k$-有理点，$N$ 是否仍总为满射。
\end{proof}

对于某些曲面，Saito 的方法使我们可以通过闭链映射控制群 $\CH^2(X)_{\tors}=\CH_0(X)_{\tors}$；更引人注意的是，还可以通过\term{Brauer 群}{Brauer group}{groupe de Brauer}来控制它。下述定理尤其适用于满足 $H^2(X,\mathcal O_X)=0$ 和 $H^1(X,\mathcal O_X)=0$ 的曲面（后一个条件保证 Albanese 簇平凡），特别是 Enriques 曲面和有理曲面。此前只知道有理曲面的情形（[C1]）。

\begin{statement}{定理 8.4（Saito，[S2]）}
设 $k$ 为局部域，$X$ 为 $k$ 上的光滑、射影、几何整曲面。假设 $H^2(X,\mathcal O_X)=0$，且 $X$ 的 Albanese 簇 $\Alb_X$ 具有潜在好约化。\footnote{译注：原书所附 2008 年勘误特别指出，此处 Albanese 簇的好约化假设不能删去；R.~Parimala 与 V.~Suresh 的例子表明了这一点，见 \textit{Inventiones math.} \textbf{122} (1995), 83--117。参见书末\hyperref[errata-ct]{《勘误与说明》}。}则：
\begin{enumerate}[label=(\alph*)]
\item 存在整数 $N>0$，使对每个 $N$ 的整数倍 $n>0$，闭链映射
\[
 \CH_0(X)_{\tors}\longrightarrow H^4_{\et}(X,\mu_n^{\otimes2})
\]
为单射。
\item 自然配对
\[
 \CH_0(X)\times\Br(X)\longrightarrow\Br(k)=\mathbb Q/\mathbb Z
\]
诱导单射
\[
 \CH_0(X)_{\tors}\hookrightarrow\Hom(\Br(X),\mathbb Q/\mathbb Z).
\]
\end{enumerate}
\end{statement}
\begin{proof}
由定理 7.1，要证明 (a)，只须证明映射
\[
 H^1(X,\mathcal K_2)\otimes\mathbb Q/\mathbb Z
 \longrightarrow H^3_{\et}(X,\mathbb Q/\mathbb Z(2))
\]
的余核指数有界。为此，重用定理 7.3 的证明。通过 Hochschild--Serre 谱序列分析 $H^3_{\et}(X,\mathbb Q/\mathbb Z(2))$。群 $H^3(k,\mathbb Q/\mathbb Z(2))$ 为零，因为局部域 $k$ 的上同调维数为 $\operatorname{cd}(k)=2$。$X$ 为局部域上的曲面，所以 $H^3_{\et}(\overline X,\mathbb Q/\mathbb Z(2))^G$ 有限（见上面 8.2(d) 的证明）。仿照 7.3 可知，这里只须证明
\origpage{46}{41}
$H^2(k,H^1_{\et}(\overline X,\mathbb Q/\mathbb Z(2)))$ 指数有界。在 7.3 中，$H^1(X,\mathcal O_X)=0$ 保证 $H^1_{\et}(\overline X,\mathbb Q/\mathbb Z(2))$ 的最大可除子群为零，因而很容易知道 $H^2(k,H^1_{\et}(\overline X,\mathbb Q/\mathbb Z(2)))$ 指数有界。

而在这里，群 $H^2(k,H^1_{\et}(\overline X,\mathbb Q/\mathbb Z(2)))$ 是有限的：Albanese 簇（从而 Picard 簇）在 $k$ 的某个有限扩张 $K$ 上具有好约化的假设，保证群
\[
 H^2(K,H^1_{\et}(\overline X,\mathbb Q/\mathbb Z(2))^0)
\]
为零\footnote{译注：此式中原印本的首个参数误作 $k$；已按 2008 年勘误改为 $K$。}（参见 [CR1]，第 190--191 页），从而群 $H^2(k,H^1_{\et}(\overline X,\mathbb Q/\mathbb Z(2))^0)$ 有限。实际上它为零，例如由 $\operatorname{cd}(k)\le2$ 可知该群可除。于是 (a) 得证。

再证 (b)。将局部类域论的对偶性与可分闭域上光滑射影簇的 Poincaré 对偶性结合，可为曲面 $X$ 得到有限阿贝尔群之间的非退化配对
\[
 H^4_{\et}(X,\mu_n^{\otimes2})\times H^2_{\et}(X,\mu_n)
 \longrightarrow H^6_{\et}(X,\mu_n^{\otimes3})=\mathbb Z/n.
\]
见 [S1]，2.9；有限域上的类似对偶见 [CSS]，第 790--791 页，(i) 的证明。Kummer 正合列则给出
\[
 0\longrightarrow\Pic(X)/n\longrightarrow H^2_{\et}(X,\mu_n)
 \longrightarrow{}_n\Br(X)\longrightarrow0.
\]
可以验证，由上述配对、闭链映射及 Kummer 正合列所得映射诱导的配对
\[
 \CH^2(X)/n\times\Pic(X)/n
 \longrightarrow H^6_{\et}(X,\mu_n^{\otimes3})=\mathbb Z/n
\]
为零：闭链映射把闭点 $P$ 映入支集在该点的上同调，而 $X$ 上每个除子类都有支集避开 $P$ 的代表。取满足 (a) 的 $N$ 的倍数 $n$，便由 (a) 得到单射
\[
 \CH_0(X)_{\tors}\hookrightarrow\Hom({}_n\Br(X),\mathbb Z/n),
\]
进而得到所述单射 $\CH_0(X)_{\tors}\hookrightarrow\Hom(\Br(X),\mathbb Q/\mathbb Z)$。稍作计算可验证，这个映射确实来自自然配对
\[
 \CH_0(X)\times\Br(X)\longrightarrow\Br(k)=\mathbb Q/\mathbb Z.
\]
类似计算见 [CSS]，第 791--792 页，(ii) 的证明。
\end{proof}

\noindent\textbf{备注 8.4.1。} 若要去掉这个不太令人满意的好约化假设，只须证明 $H^3_{\et}(X,\mathbb Q/\mathbb Z(2))$ 在 $H^3_{\et}(k(X),\mathbb Q/\mathbb Z(2))$ 中的像有限，或至少指数有界。事实上，无须假设曲面 $X$ 具有好约化，已知 $\CH^2(X)_{\tors}$ 有限（定理 8.2）。若上述像有限，图 (7.2) 说明群 $\Theta$ 指数有界，因而可应用定理 7.1。\footnote{译注：本备注保留原文提出问题时的表述。2008 年勘误说明，定理 8.4 中 Albanese 簇的好约化假设不能去掉；见上一定理脚注及书末\hyperref[errata-ct]{《勘误与说明》}。}

\noindent\textbf{备注 8.4.2。} 在 8.4 的假设下，另加一个大概是多余的假设（即 Bloch 关于满足 $H^2(X,\mathcal O_X)=0$ 的曲面的猜想）：Albanese 映射 $A_0(\overline X)\to\Alb_X(\overline k)$ 为同构。Saito [S2] 证明，此时配对
\[
 A_0(X)\times\Br(X)\longrightarrow\mathbb Q/\mathbb Z
\]
实际上左非退化。其思想是通过局部域上阿贝尔簇的 Tate 对偶来控制 $A_0(X)$ 在 $\Alb_X(k)$ 中的像。Tate 对偶给出同构
\[
 \Alb_X(k)\simeq\Hom(H^1(k,\underline{\Pic}^{\,0}_X),\mathbb Q/\mathbb Z).
\]

\origpage{47}{42}
自然复合同态\translatornote{3}
\[
 H^1(k,\underline{\Pic}^{\,0}_X)
 \longrightarrow H^1(k,\underline{\Pic}_X)\longrightarrow\Br(X)
\]
给出下图，可以验证其交换性：
\[
\begin{tikzcd}[ampersand replacement=\&,column sep=large,row sep=large]
 A_0(X)\arrow[r]\arrow[d]\&\Alb_X(k)\arrow[d,"{\simeq}"]\\
 \Hom(\Br(X),\mathbb Q/\mathbb Z)\arrow[r]
 \&\Hom(H^1(k,\underline{\Pic}^{\,0}_X),\mathbb Q/\mathbb Z).
\end{tikzcd}
\]
因此左侧竖直箭头的核包含在上方水平箭头的核中；由附加假设，后者为挠群。上一定理又保证，左侧竖直箭头在 $A_0(X)_{\tors}$ 上的限制为单射。

\begin{statement}{定理 8.5}
设 $X$ 是局部域 $k$ 上的光滑、射影、几何整曲面。假设：
\begin{enumerate}[label=(\roman*)]
\item $H^2(X,\mathcal O_X)=0$；
\item Albanese 映射 $A_0(\overline X)\to\Alb_X(\overline k)$ 为同构（按照 Bloch 的一个猜想，这由 (i) 推出）。
\end{enumerate}
则 $A_0(X)$ 是 $\Alb_X(k)$ 的一个开子群被有限群的扩张。特别地，对每个整数 $n>0$，商群 $A_0(X)/n$ 有限；对几乎所有素数 $l$，商群 $A_0(X)/l$ 为零。
\end{statement}
\begin{proof}
由 (ii)，$A_0(X)\to\Alb_X(k)$ 的核为挠群，所以在 (i) 的假设下，由定理 8.2(d) 可知其有限。下面研究 $A_0(X)$ 在 $\Alb_X(k)$ 中的像。后者是交换的紧 $p$-进解析群，其开子群恰为有限指数子群。设 $C$ 是 $X$ 的一条光滑超平面截面曲线。有阿贝尔簇的满同态 $\Alb_C\to\Alb_X$，它诱导的同态 $\Alb_C(k)\to\Alb_X(k)$ 的像为开子群，故在 $\Alb_X(k)$ 中具有有限指数。若 $C(k)\ne\varnothing$，映射 $A_0(C)\to\Alb_C(k)$ 为同构；一般情形下，它的像是 $\Alb_C(k)$ 的有限指数子群（因为 $k$ 为局部域）。所以 $A_0(X)\to\Alb_X(k)$ 的像包含一个有限指数子群，其自身因而也具有有限指数，从而是开子群。紧 $p$-进解析群 $\Alb_X(k)$ 是一个有限群被同构于有限个 $\mathbb Z_p$ 的直和的子群的扩张；这就推出结论的最后一部分。
\end{proof}

最后，不加证明地引述一个用局部化方法得到的结果。
\begin{statement}{定理 8.6（[CR3]，6.3）}
设 $k$ 为局部域，即 $\mathbb Q_p$ 的有限扩张，$R$ 为其整数环，$\mathbb F$ 为其剩余域。设 $\mathcal X$ 为 $R$ 上的光滑射影概形，几何纤维均为整概形，且相对维数为 $2$；\footnote{译注：“相对维数为 $2$”这一必要假设根据原书所附 2008 年勘误补入。}其一般纤维为 $X/k$，特殊纤维为 $Y/\mathbb F$。假设 $H^2(Y,\mathcal O_Y)=0$。则对每个素数 $l\ne p$，$l$-准素挠子群 $\CH^2(X)_{l\text{-}\tors}$ 与 $\CH^2(Y)_{l\text{-}\tors}$ 都有限，并且自然同构。
\end{statement}

\origpage{48}{43}
\section*{§9. 数域上的簇，III}
\addcontentsline{toc}{section}{§9. 数域上的簇，III}
前面几节证明了：当 $X$ 满足 $H^2(X,\mathcal O_X)=0$ 时，$\CH^2(X)$ 的挠元的指数有界。至少在整体情形，这些证明主要归功于 Salberger（整体域上的定理 6.6；局部域上的定理 8.2(b)，后者带有一些附加假设）。本节首先介绍 Salberger 最初证明挠元指数有界的方法（定理 9.1 及定理 9.2 证明的开头）。这里的证明是根据 Salberger 1990 年的报告整理出来的。

随后介绍 Salberger 如何修改 §6 中的局部化技巧（私人通信），从而证明有限性定理 6.1，即此处的定理 9.2，而无须像 Raskind 和我那样，使用 Bloch、Kato--Saito 和 Somekawa 的精细结果（定理 6.5 证明中的 c)）。

最近的预印本 [Sb2] 还包含许多其他有趣的思想和证明，建议读者参阅。
\begin{statement}{定理 9.1}
设 $k$ 是特征零的非阿基米德局部域或整体域，$X$ 是 $k$ 上的光滑、射影、几何整簇。假设 $H^2(X,\mathcal O_X)=0$。则映射
\[
\begin{split}
 H^1(X,\mathcal K_2)\otimes\mathbb Q/\mathbb Z
 \longrightarrow\Ker\bigl[H^3_{\et}(X,\mathbb Q/\mathbb Z(2))
 \longrightarrow {}&H^3_{\et}(\overline X,\mathbb Q/\mathbb Z(2))^G\\
 &\oplus H^3_{\et}(k(X),\mathbb Q/\mathbb Z(2))\bigr]
\end{split}
\]
的余核指数有界。
\end{statement}
\begin{proof}
仿照定理 7.3 的证明。平展上同调的 Hochschild--Serre 谱序列
\[
 H^p\bigl(k,H^q_{\et}(\overline X,\mathbb Q/\mathbb Z(2))\bigr)
 \Longrightarrow H^{p+q}_{\et}(X,\mathbb Q/\mathbb Z(2))
\]
在群 $\Ker[H^3_{\et}(X,\mathbb Q/\mathbb Z(2))\to H^3_{\et}(\overline X,\mathbb Q/\mathbb Z(2))^G]$ 上给出过滤 $F_X^0\subset F_X^1\subset F_X^2$。记
\[
\begin{split}
 N_rH^3_{\et}(X,\mathbb Q/\mathbb Z(2))
 =\Ker\bigl[H^3_{\et}(X,\mathbb Q/\mathbb Z(2))
 \longrightarrow {}&H^3_{\et}(\overline X,\mathbb Q/\mathbb Z(2))^G\\
 &\oplus H^3_{\et}(k(X),\mathbb Q/\mathbb Z(2))\bigr],
\end{split}
\]
并考虑此子群上的诱导过滤
\[
 F_{X,r}^0\subset F_{X,r}^1\subset F_{X,r}^2.
\]
群 $F_{X,r}^0$ 是 $F_X^0$ 的子群，而 $F_X^0$ 又是 $H^3(k,\mathbb Q/\mathbb Z(2))$ 的商群。在局部域上，后一个群有限；在数域上，它等于 $(\mathbb Z/2)^s$，其中 $s$ 是 $k$ 的实完备化的个数。事实上，在任意域 $k$ 上，若 $X(k)\ne\varnothing$，则由特殊化论证可知 $F_{X,r}^0=0$。

同定理 7.3 中一样，条件 $H^2(X,\mathcal O_X)=0$ 蕴含映射
\[
 H^1(X,\mathcal K_2)\otimes\mathbb Q/\mathbb Z\longrightarrow F_X^2/F_X^1
\]
的余核指数有界；因此映射
\[
 H^1(X,\mathcal K_2)\otimes\mathbb Q/\mathbb Z
 \longrightarrow F_{X,r}^2/F_{X,r}^1
\]
的余核也指数有界。

故只须证明 $F_{X,r}^1$ 指数有界。为此，再使用 6.7 和 8.3 中用过的超平面截面技巧。设 $C\subset X$ 是通过逐次取超平面截面得到的光滑、射影、几何整 $k$-曲线。

\origpage{49}{44}
Hochschild--Serre 谱序列对任意态射都具有反变函子性，因而得到交换图
\[
\begin{tikzcd}[ampersand replacement=\&,column sep=large,row sep=large]
 0\arrow[r]\&F_X^0\arrow[r]\arrow[d]\&F_X^1\arrow[r]\arrow[d]
 \&E_{\infty,X}^{21}\arrow[r]\arrow[d]\&0\\
 0\arrow[r]\&F_C^0\arrow[r]\&F_C^1\arrow[r]\&E_{\infty,C}^{21}\arrow[r]\&0,
\end{tikzcd}
\]
其中 $F_X^0$ 与 $F_C^0$ 都有限。另外，对局部域或整体域 $k$，$H^4(k,\mathbb Q/\mathbb Z(2))=0$。所以，利用谱序列的函子性，得到正合列交换图\translatornoteref{4}
\[
\resizebox{.98\linewidth}{!}{$
\begin{tikzcd}[ampersand replacement=\&,column sep=large,row sep=large]
 \&H^2(k,H^1_{\et}(\overline X,\mathbb Q/\mathbb Z(2)))\arrow[r]\arrow[d]
 \&E_{\infty,X}^{21}\arrow[r]\arrow[d]\&0\\
 H^0(\overline C,\mathbb Q/\mathbb Z(2))^G\arrow[r]
 \&H^2(k,H^1_{\et}(\overline C,\mathbb Q/\mathbb Z(2)))\arrow[r]
 \&E_{\infty,C}^{21}\arrow[r]\&0.
\end{tikzcd}$}
\]
图中的群 $H^0(\overline C,\mathbb Q/\mathbb Z(2))^G=\mathbb Q/\mathbb Z(1)^G=\mu(k)$ 是 $k$ 中的单位根群，因而有限。\translatornote{4}6.7（整体情形）及 8.3（局部情形）中的论证表明，左侧竖直箭头的核指数有界。图追踪便说明，$F_X^1\to F_C^1$ 的核也指数有界。

设 $Y\subset X$ 是光滑且几何整的超平面截面，$R$ 是 $X$ 在 $Y$ 的一般点处的局部环。由 Bloch--Ogus 理论 [Bl--O]，$NH^3_{\et}(X,\mathbb Q/\mathbb Z(2))$ 中任意元素在 $H^3_{\et}(R,\mathbb Q/\mathbb Z(2))$ 中的像都为零，因此传到 $R$ 的剩余域 $k(Y)$ 后，其在 $H^3_{\et}(k(Y),\mathbb Q/\mathbb Z(2))$ 中的像也为零。归纳可知，限制映射
\[
 H^3_{\et}(X,\mathbb Q/\mathbb Z(2))\longrightarrow H^3_{\et}(C,\mathbb Q/\mathbb Z(2))
\]
把子群 $NH^3_{\et}(X,\mathbb Q/\mathbb Z(2))$ 映入 $NH^3_{\et}(C,\mathbb Q/\mathbb Z(2))$，也把 $N_rH^3_{\et}(X,\mathbb Q/\mathbb Z(2))$ 映入 $N_rH^3_{\et}(C,\mathbb Q/\mathbb Z(2))$。

特别地，$F_{X,r}^1\to F_C^1$ 的像包含在
\[
 F_{C,r}^1=\Ker[F_C^1\longrightarrow H^3(k(C),\mathbb Q/\mathbb Z(2))]
\]
中。比较 $C$ 与函数域 $k(C)$ 的 Hochschild--Serre 谱序列，得到交换图
\[
\resizebox{.98\linewidth}{!}{$
\begin{tikzcd}[ampersand replacement=\&,column sep=large,row sep=large]
 0\arrow[r]\&F_C^0\arrow[r]\arrow[d]\&F_C^1\arrow[r]\arrow[d]
 \&H^2(k,H^1_{\et}(\overline C,\mathbb Q/\mathbb Z(2)))/\im(\mu(k))\arrow[r]\arrow[d]\&0\\
 0\arrow[r]\&F_{k(C)}^0\arrow[r]\&F_{k(C)}^1\arrow[r]
 \&H^2(k,H^1(\overline k(C),\mathbb Q/\mathbb Z(2)))\arrow[r]\&0,
\end{tikzcd}$}
\]
其中 $F_C^0$ 与 $F_{k(C)}^0$ 都是 $H^3(k,\mathbb Q/\mathbb Z(2))$ 的有限商群。所以，忽略有限群后，$F_{C,r}^1$ 等同于
\[
 \Ker[H^2(k,H^1_{\et}(\overline C,\mathbb Q/\mathbb Z(2)))
 \longrightarrow H^2(k,H^1(\overline k(C),\mathbb Q/\mathbb Z(2)))].
\]
而这个群正是
\[
 \Ker[H^2(G,H^0(\overline C,\mathcal K_2))
 \longrightarrow H^2(G,K_2\overline k(C))],
\]
\origpage{50}{45}
我们在 6.7 和 8.3 中已说明：对于局部域或整体域上的曲线，该群指数有界；若曲线有有理点，则该群为零（[R3]）。

由于 $F_{X,r}^1\to F_{C,r}^1$ 的核指数有界，可知 $F_{X,r}^1$ 也指数有界，证明完成。

容易验证，在局部情形，上述证明中所有“指数有界”都可以换成“有限”。
\end{proof}

\begin{statement}{定理 9.2（即定理 6.1）}
设 $X$ 是数域 $k$ 上的光滑、射影、几何连通簇。假设 $H^2(X,\mathcal O_X)=0$，则挠子群 $\CH^2(X)_{\tors}$ 有限。
\end{statement}
\begin{proof}
先证明 $\CH^2(X)_{\tors}$ 指数有界。比较正合列
\[
 0\longrightarrow H^1(X,\mathcal K_2)\otimes\mathbb Q/\mathbb Z
 \longrightarrow NH^3_{\et}(X,\mathbb Q/\mathbb Z(2))
 \longrightarrow\CH^2(X)_{\tors}\longrightarrow0
\]
及其在代数闭包上的对应正合列
\[
 0\longrightarrow H^1(\overline X,\mathcal K_2)\otimes\mathbb Q/\mathbb Z
 \longrightarrow NH^3_{\et}(\overline X,\mathbb Q/\mathbb Z(2))
 \longrightarrow\CH^2(\overline X)_{\tors}\longrightarrow0.
\]
利用 $H^1(\overline X,\mathcal K_2)\otimes\mathbb Q/\mathbb Z=0$（[CR1]）及 $H^3_{\et}(\overline X,\mathbb Q/\mathbb Z(2))^G$ 的有限性（参见定理 6.4 的证明），由定理 9.1 可得 $\CH^2(X)_{\tors}$ 指数有界。

于是，对几乎所有素数 $l$，$\CH^2(X)_{l\text{-}\tors}=0$；要证明 $\CH^2(X)_{\tors}$ 有限，只须对任意给定素数 $l$ 证明 $\CH^2(X)_{l\text{-}\tors}$ 有限。

现在回到定理 6.5 的证明，沿用那里的记号。参看其中 b)，有局部化正合列
\[
 H^1(X_L,\mathcal K_2)\longrightarrow
 \bigoplus_{q\in\Spec(B)^{(1)}}\Pic(X_q)
 \longrightarrow\CH^2(\mathcal X_B)
 \longrightarrow\CH^2(X_L)\longrightarrow0,
\]
而且已经证明，复合映射
\[
 H^1(X_L,\mathcal K_2)\longrightarrow
 \bigoplus_{q\in\Spec(B)^{(1)}}\Pic(X_q)
 \longrightarrow\bigoplus_{q\in\Spec(B)^{(1)}}\NS(X_q)
\]
为满射。因此，诱导的复合映射
\[
 H^1(X_L,\mathcal K_2)\otimes\mathbb Q\longrightarrow
 \bigoplus_{q\in\Spec(B)^{(1)}}\Pic(X_q)\otimes\mathbb Q
 \longrightarrow\bigoplus_{q\in\Spec(B)^{(1)}}\NS(X_q)\otimes\mathbb Q
\]
也为满射。右侧箭头在这里是同构，因为群
\[
 J(\mathbb F_q)=\Ker[\Pic(X_q)\longrightarrow\NS(X_q)]
\]
都是挠群，甚至是有限群，所以 $J(\mathbb F_q)\otimes\mathbb Q=0$。因此
\[
 H^1(X_L,\mathcal K_2)\otimes\mathbb Q\longrightarrow
 \bigoplus_{q\in\Spec(B)^{(1)}}\Pic(X_q)\otimes\mathbb Q
\]
是满射。利用从 $L$ 到 $k$ 的迹，又可知
\[
 H^1(X,\mathcal K_2)\otimes\mathbb Q\longrightarrow
 \bigoplus_{p\in\Spec(A)^{(1)}}\Pic(X_p)\otimes\mathbb Q
\]
为满射。\footnote{译注：这一使用迹将满射性降回 $k$ 的步骤，按照 2008 年勘误补入。}
由此满射性及上述局部化正合列，形式上的论证表明：Chow 群到一般纤维的限制映射当然是满射，而且它在挠子群上也诱导满射
\origpage{51}{46}
\[
 \CH^2(\mathcal X)_{\tors}\longrightarrow\CH^2(X)_{\tors}.
\]
\footnote{译注：原印本这里写作 $\CH^2(\mathcal X_B)_{\tors}\to\CH^2(X_L)_{\tors}$；按 2008 年勘误，应用迹后应回到 $A/k$ 上的模型与一般纤维。勘误页左端仍印下标 $2$；此处依余维二 Chow 群的上下文写为上标 $2$。}

现在置于定理 6.2 的情形，并沿用其中的记号。让整数 $n$ 遍历素数 $l$ 的幂，得到下图：
\[
\begin{tikzcd}[ampersand replacement=\&,column sep=large,row sep=large]
 H^1(\mathcal X,\mathcal H^2(\mathbb Q_l/\mathbb Z_l(2)))\arrow[r]\arrow[d,"{\gamma_2}"']
 \&\CH^2(\mathcal X)_{l\text{-}\tors}\arrow[r]\&0\\
 H^3_{\et}(\mathcal X,\mathbb Q_l/\mathbb Z_l(2)).
\end{tikzcd}
\]
其中 $\gamma_2$ 为单射，水平箭头为满射，因为滤过归纳极限保持满射与单射。

$l$-准素挠群 $H^3_{\et}(\mathcal X,\mathbb Q_l/\mathbb Z_l(2))$ 是有限群 $H^3_{\et}(\mathcal X,\mu_{l^n}^{\otimes2})$ 的正向极限，它具有\term{有限余型}{finite cotype}{cotype fini}，即它是一个有限 $l$-准素群被某个 $(\mathbb Q_l/\mathbb Z_l)^m$ 的扩张；也就是说，它是一个有限 $l$-准素群与有限个 $\mathbb Q_l/\mathbb Z_l$ 的直和。更一般地，对每个 $H^i_{\et}(\mathcal X,\mathbb Q_l/\mathbb Z_l(j))$ 都有这一结论。为此，采用下述经典论证（参见 [M1]，第 163--166 页，及 [CSS]，第 773--774、780--781 页）。考虑由 $\mathcal X$ 上的平展层正合列
\[
 1\longrightarrow\mu_{l^n}^{\otimes j}
 \longrightarrow\mu_{l^{n+m}}^{\otimes j}
 \longrightarrow\mu_{l^m}^{\otimes j}\longrightarrow1
\]
得到的上同调长正合列。先对 $n$ 取投射极限；由于群 $H^i_{\et}(\mathcal X,\mu_{l^n}^{\otimes j})$ 都有限（见定理 6.2 证明中的 b)），所得列仍正合。再对所得长正合列中的 $m$ 取归纳极限，便得到短正合列
\[
\begin{split}
 0\longrightarrow
 H^i(\mathcal X,\mathbb Z_l(j))/H^i(\mathcal X,\mathbb Z_l(j))_{\tors}
 &\longrightarrow H^i(\mathcal X,\mathbb Q_l(j))\\
 &\longrightarrow H^i(\mathcal X,\mathbb Q_l/\mathbb Z_l(j))\\
 &\longrightarrow H^{i+1}(\mathcal X,\mathbb Z_l(j))_{l\text{-}\tors}
 \longrightarrow0,
\end{split}
\]
其中 $H^i(\mathcal X,\mathbb Q_l(j))=H^i(\mathcal X,\mathbb Z_l(j))\otimes_{\mathbb Z_l}\mathbb Q_l$。\footnote{译注：短正合列首个非零项按 2008 年勘误改为 $H^i(\mathcal X,\mathbb Z_l(j))$ 对其挠子群的商群。}
注意，对任意整数 $i,j$，群 $H^i_{\et}(\mathcal X,\mu_{l^n}^{\otimes j})$ 的投射极限都是有限生成的 $\mathbb Z_l$-模。由上述正合列，确可推出 $H^i(\mathcal X,\mathbb Q_l/\mathbb Z_l(j))$ 是有限 $l$-准素群被若干 $\mathbb Q_l/\mathbb Z_l$ 的直和的扩张，因而实际上是一个有限 $l$-准素群与若干 $\mathbb Q_l/\mathbb Z_l$ 的直和。

利用紧 $\mathbb Z_l$-模与离散 $l$-准素挠模之间的对偶性可知，这一性质传递到任何子商群。因此，$\CH^2(\mathcal X)_{l\text{-}\tors}$ 作为 $H^3_{\et}(\mathcal X,\mathbb Q_l/\mathbb Z_l(2))$ 的子商群，具有有限余型；$\CH^2(X)_{l\text{-}\tors}$ 也如此，因为已证明它是前者的商群。

然而，证明开头已知 $\CH^2(X)_{l\text{-}\tors}$ 指数有界。因此它是有限群。
\end{proof}

\origpage{52}{47}
\section*{参考文献}
\addcontentsline{toc}{section}{第一篇参考文献}
\begin{description}[leftmargin=3.4em,labelsep=.5em,font=\normalfont]
\item[{[BV]}] L.~Barbieri-Viale. Cohomology theories and algebraic cycles on singular varieties, \textit{Contemp. Math.} \textbf{126} (1992), 197--217; Theorie coomologiche e cicli algebrici sulle varietà singolari, Tesi di dottorato, Università di Genova (1991).
\item[{[B1]}] S.~Bloch. Torsion algebraic cycles and a theorem of Roitman, \textit{Comp. Math.} \textbf{39} (1979), 107--127.
\item[{[B2]}] S.~Bloch. \textit{Lectures on algebraic cycles}, Duke Univ. Math. Ser.~4, Durham, N.C., 1980.
\item[{[B3]}] S.~Bloch. Torsion algebraic cycles, $K_2$ and Brauer groups of function fields, 载 L.N.M.~\textbf{844}（M.~Kervaire、M.~Ojanguren 编）, Springer, 1981.
\item[{[B4]}] S.~Bloch. On the Chow groups of certain rational surfaces, \textit{Ann. Sc. Éc. Norm. Sup.} \textbf{14} (1981), 41--59.
\item[{[B5]}] S.~Bloch. Algebraic $K$-theory and class field theory for arithmetic surfaces, \textit{Ann. of Math.} \textbf{114} (1981), 229--265.
\item[{[BO]}] S.~Bloch、A.~Ogus. Gersten's conjecture and the homology of schemes, \textit{Ann. Sc. E.N.S.} \textbf{7} (1974), 181--202.
\item[{[BS]}] S.~Bloch、V.~Srinivas. Remarks on correspondences and algebraic cycles, \textit{Am. J. Math.} \textbf{105} (1983), 1235--1253.
\item[{[C1]}] J.-L.~Colliot-Thélène. Hilbert's theorem 90 for $K_2$, with application to the Chow group of rational surfaces, \textit{Invent. Math.} \textbf{35} (1983), 1--20.
\item[{[C2]}] J.-L.~Colliot-Thélène. On the reciprocity sequence in higher class field theory of function fields, Prépublication d'Orsay \textbf{92-52} (1992).
\item[{[CR1]}] J.-L.~Colliot-Thélène、W.~Raskind. $K_2$-cohomology and the second Chow group, \textit{Math. Ann.} \textbf{270} (1985), 165--199.
\item[{[CR2]}] J.-L.~Colliot-Thélène、W.~Raskind. On the reciprocity law for surfaces over finite fields, \textit{J. Fac. Sc. Univ. Tokyo, Sect. IA} \textbf{33} (1986), 283--294.
\item[{[CR3]}] J.-L.~Colliot-Thélène、W.~Raskind. Groupe de Chow de codimension deux des variétés définies sur un corps de nombres : un théorème de finitude pour la torsion, \textit{Invent. Math.} \textbf{105} (1991), 221--245.
\item[{[CS]}] J.-L.~Colliot-Thélène、J.-J.~Sansuc. On the Chow groups of certain rational surfaces : a sequel to a paper of S.~Bloch, \textit{Duke Math. J.} \textbf{48} (1981), 421--447.
\item[{[CSS]}] J.-L.~Colliot-Thélène、J.-J.~Sansuc、C.~Soulé. Torsion dans le groupe de Chow de codimension deux, \textit{Duke Math. J.} \textbf{50} (1983), 763--801.
\item[{[Cb]}] K.~Coombes. The arithmetic of zero-cycles on surfaces with geometric genus and irregularity zero, \textit{Math. Ann.} \textbf{291} (1991), 429--452.
\item[{[F]}] W.~Fulton. \textit{Intersection Theory}, Ergeb. der Math. und ihr. Grenzgeb., 3.~Folge, Bd.~2, Springer, 1984.
\item[{[G]}] M.~Gros. $0$-cycles de degré zéro sur les surfaces fibrées en coniques, \textit{J. reine und angew. Math.} \textbf{373} (1987), 166--184.
\end{description}

\origpage{53}{48}
\begin{description}[leftmargin=3.4em,labelsep=.5em,font=\normalfont]
\item[{[GS]}] M.~Gros、N.~Suwa. Application d'Abel-Jacobi $p$-adique et cycles algébriques, \textit{Duke Math. J.} \textbf{57} (1988), 579--613.
\item[{[J]}] U.~Jannsen. On the $l$-adic cohomology of varieties over number fields and its Galois cohomology, 载 \textit{Galois groups over $\mathbb Q$}（Y.~Ihara、K.~Ribet、J.-P.~Serre 编）, Springer-Verlag, 1987.
\item[{[K]}] K.~Kato. A Hasse principle for two dimensional global fields, \textit{J. für die reine und angew. Math.} \textbf{366} (1986), 142--181.
\item[{[KS]}] K.~Kato、S.~Saito. Unramified class field theory of arithmetical surfaces, \textit{Ann. of Math.} \textbf{118} (1983), 241--275.
\item[{[L]}] F.~Lecomte. Rigidité des groupes de Chow, \textit{Duke Math. J.} \textbf{53} (1986), 405--426.
\item[{[MS]}] A.~S.~Merkur'ev、A.~A.~Suslin. $K$-cohomology of Severi-Brauer varieties and norm residue homomorphism, \textit{Izv. Akad. Nauk SSSR} \textbf{46} (1982), 1011--1146 = \textit{Math. USSR Izv.} \textbf{21} (1983), 307--341.
\item[{[M1]}] J.~S.~Milne. \textit{Etale cohomology}, Princeton University Press, 33, Princeton, N.J., 1980.
\item[{[M2]}] J.~S.~Milne. \textit{Arithmetic Duality Theorems}, Perspect. Math.~1, Academic Press, Boston, 1986.
\item[{[Q]}] D.~Quillen. Higher algebraic $K$-theory I, 载 \textit{Algebraic $K$-theory I}（H.~Bass 编）, LNM~\textbf{341}, Springer-Verlag, 1973.
\item[{[R1]}] W.~Raskind. Algebraic $K$-theory, étale cohomology and torsion algebraic cycles, 载 \textit{Algebraic $K$-theory and Algebraic Number Theory}, Honolulu 1987, Contemporary Mathematics \textbf{83} (1989), 311--341.
\item[{[R2]}] W.~Raskind. Torsion algebraic cycles on varieties over local fields, 载 \textit{Algebraic $K$-theory : Connections with Geometry and Topology}, Lake Louise 1987（J.~F.~Jardine、V.~P.~Snaith 编）, Kluwer Academic Publishers, 1989.
\item[{[R3]}] W.~Raskind. On $K_1$ of curves over global fields, \textit{Math. Ann.} \textbf{288} (1990), 179--193.
\item[{[S1]}] S.~Saito. A global duality theorem for varieties over global fields, 载 \textit{Algebraic $K$-theory : Connections with Geometry and Topology}, Lake Louise 1987（J.~F.~Jardine、V.~P.~Snaith 编）, Kluwer Academic Publishers, 1989.
\item[{[S2]}] S.~Saito. A conjecture of Bloch and Brauer groups of surfaces over $p$-adic fields, 预印本，1990 年 5 月。
\item[{[S3]}] S.~Saito. Torsion zero-cycles and étale homology of singular schemes, \textit{Duke Math. J.} \textbf{64} (1991), 71--83.
\item[{[S4]}] S.~Saito. Cycle map on torsion algebraic cycles of codimension two, \textit{Invent. Math.} \textbf{106} (1991), 443--460.
\item[{[Sb1]}] P.~Salberger. Zero-cycles on rational surfaces over number fields, \textit{Invent. Math.} \textbf{91} (1988), 505--524.
\item[{[Sb2]}] P.~Salberger. Chow groups of codimension two and $l$-adic realizations of motivic cohomology, 打字稿，1991 年 9 月。\footnote{原书所附 2008 年勘误补充：此文已以题名 \textit{Torsion cycles of codimension 2 and $l$-adic realizations of motivic cohomology} 发表，载 \textit{Séminaire de théorie des nombres, Paris 1991--1992}（Sinnou David 编），Progress in Mathematics \textbf{116}, Birkhäuser (1993), 247--277。}
\item[{[Sc]}] C.~Schoen. Some examples of torsion in the Griffiths group, \textit{Math. Ann.} \textbf{293} (1992), 651--679.
\item[{[So]}] M.~Somekawa. On Milnor $K$-groups attached to semi-abelian varieties, \textit{Journal of $K$-theory} \textbf{4} (1990), 105--119.
\item[{[S]}] C.~Soulé. $K_2$ et le groupe de Brauer [d'après A.~S.~Merkurjev et A.~A.~Suslin], Séminaire Bourbaki, 1982/83, n\textsuperscript{o}~601.
\end{description}

\origpage{54}{49}
\begin{description}[leftmargin=3.4em,labelsep=.5em,font=\normalfont]
\item[{[Sr]}] V.~Srinivas. Rational equivalence of $0$-cycles on normal varieties over $\mathbb C$, \textit{Proceedings of Symposia in Pure Mathematics} \textbf{46}, Part~2 (1987), 475--482.
\item[{[Su]}] A.~A.~Suslin. Torsion in $K_2$ of fields, \textit{Journal of $K$-theory} \textbf{1} (1987), 5--29.
\item[{[Sw]}] N.~Suwa. Sur l'image de l'application d'Abel-Jacobi de Bloch, \textit{Bull. Soc. Math. France} \textbf{116} (1988), 69--101.
\item[{[SGA4 1/2]}] \textit{Cohomologie étale}, Séminaire de Géométrie Algébrique du Bois-Marie SGA~$4\frac12$, P.~Deligne 著，Springer L.N.M.~\textbf{569} (1977).
\end{description}
